\documentclass[utf8,twocolumn]{ctexart}
\usepackage{ctex}
\usepackage{amsmath}
\usepackage{circuitikz}
\usepackage{tikz}
\usepackage{graphicx}
\usepackage{hyperref}
\usepackage{geometry}
\usepackage{pdfpages}
\usepackage{booktabs}
\usepackage{subfigure}
\usepackage{float}
\usepackage{amsmath}
\usepackage{amssymb}
% math font
\usepackage{mathrsfs}
\usepackage{multirow}
\usepackage{inputenc}
\usepackage{fancyhdr}
\usepackage{physics}

\hypersetup{
    colorlinks=true,
    linkcolor=black,
    filecolor=black,      
    urlcolor=black,
    citecolor=black,
}
\geometry{a4paper, scale=0.8}
\pagestyle{fancy}
\fancyhf{}
\lhead{半导体激光器}
\rhead{无37~董皓彧~2023010659}
\cfoot{---~~\thepage~~---}

% 双栏间隔
\setlength{\columnsep}{20pt}

\title{半导体激光器特性测量实验报告}
\author{无37~董皓彧~2023010659}
\date{\zhtoday}


\begin{document}

\maketitle
\thispagestyle{fancy}

\section{实验过程及数据分析}

\subsection{DFB 半导体激光器的 P-I 曲线}
将 DFB 激光器的输出端使用光纤直接连接到光功率计上，调整驱动电流。先以 2.0 mA 为步长粗调，在根据粗调结果在部分区间内以 0.2 mA 为步长细调，测量并记录激光器的输出功率。保持温度为 $14.0^\circ\mathrm{C}$ 不变。测量结果如表 \ref{tab:PI}. 绘制 P-I 曲线如图 \ref{fig:PI}.
\begin{table}[H]
    \centering
    \begin{tabular}{cc|cc}
        \toprule
        $I_\mathrm{d} (\text{mA})$ & $P_\mathrm{out} (\mu\text{W})$ & $I_\mathrm{d} (\text{mA})$ & $P_\mathrm{out} (\mu\text{W})$ \\
        \midrule
        0.2  & 0.04375 & 8.2  & 11.94 \\
        2.0  & 0.5766  & 8.4  & 12.74 \\
        4.0  & 2.106   & 8.6  & 16.11 \\
        6.0  & 5.480   & 8.8  & 22.86 \\
        8.0  & 11.45   & 9.0  & 37.15 \\
        10.0 & 104.8   & 9.2  & 47.53 \\
        12.0 & 250.0   & 9.4  & 61.80 \\
        14.0 & 391.8   & 9.6  & 74.30 \\
        16.0 & 538.8   & 9.8  & 93.33 \\
        18.0 & 680.6   &      &       \\
        20.0 & 827.5   &      &       \\
        22.0 & 977.1   &      &       \\
        24.0 & 1124    &      &       \\
        25.0 & 1191    &      &       \\
        \bottomrule
    \end{tabular}
    \caption{DFB 半导体激光器的 P-I 曲线数据}
    \label{tab:PI}
\end{table}
\begin{figure}[H]
    \centering
    \includegraphics[width=\linewidth]{images/P-I-curve.png}
    \vspace{-2em}
    \caption{DFB 半导体激光器的 P-I 曲线}
    \label{fig:PI}
\end{figure}
使用直线拟合法得到阈值电流为 $8.55\text{ mA}$，斜率效率为 $0.07244\text{ mW/mA}$。 


\subsection{DFB 半导体激光器的 P-T 规律}
固定驱动电流为 $10.0\text{ mA}$ 不变，调整温控电流源的温度调节旋钮，记录不同温度下的输出功率。测量结果如表 \ref{tab:PT}. 绘制 P-T 曲线如图 \ref{fig:PT}.
\begin{figure}[H]
    \centering
    \includegraphics[width=\linewidth]{images/P-T-curve.png}
    \vspace{-2em}
    \caption{DFB 半导体激光器的 P-T 曲线}
    \label{fig:PT}
\end{figure}
\begin{table}[H]
    \centering
    \begin{tabular}{cc}
        \toprule
        $T (^\circ\mathrm{C})$ & $P_\mathrm{out} (\mu\text{W})$ \\
        \midrule
        14.0 & 105.7 \\ 
        15.0 & 92.14 \\ 
        16.0 & 77.86 \\ 
        17.0 & 64.11 \\ 
        18.0 & 48.53 \\ 
        19.0 & 34.78 \\ 
        20.0 & 22.06 \\ 
        21.0 & 15.05 \\ 
        22.0 & 13.49 \\ 
        23.0 & 12.45 \\ 
        23.9 & 11.93 \\ 
        25.0 & 11.16 \\ 
        \bottomrule
    \end{tabular}
    \caption{DFB 半导体激光器的 P-T 规律数据}
    \label{tab:PT}
\end{table}

\subsection{测量 DFB 和 FP 半导体激光器光谱}
\subsubsection{DFB 激光器}
将 DFB 激光器的输出端使用光纤直接连接到光谱仪上，调整驱动电流为 $10.0\text{ mA}$，分别调整温度为 $14.0^\circ\mathrm{C}$、$16.0^\circ\mathrm{C}$ 和 $20.0^\circ\mathrm{C}$，测量得光谱如图 \ref{fig:DFB-spectrum}. 发现随着温度的升高，激光器的峰值波长变长，输出功率变小，但谱线宽度基本不变。
\subsubsection{FP 激光器}
将 FP 激光器的输出端使用光纤直接连接到光谱仪上，调整驱动电流为 $10.0\text{ mA}$，温度为 $14.0^\circ\mathrm{C}$，测量得光谱如图 \ref{fig:FP-spectrum}. 发现 FP 激光器的光谱中存在多个峰值，且峰值之间间隔较为均匀，说明这些峰值对应于激光器腔内的不同纵模。整体包络呈高斯分布形状。
\begin{figure}[H]
    \centering
    \subfigure[$14.0^\circ\mathrm{C}$，峰值：$1550.4160\text{ nm},-9.88\text{ dBm}$；3dB衰减宽度为$0.180\text{ nm}$]{\includegraphics[width=\linewidth]{images/DFB-spectrum-14.jpg}} \\
    \subfigure[$16.0^\circ\mathrm{C}$，峰值：$1550.5920\text{ nm},-11.59\text{ dBm}$；3dB衰减宽度为$0.180\text{ nm}$]{\includegraphics[width=\linewidth]{images/DFB-spectrum-16.jpg}} \\
    \subfigure[$20.0^\circ\mathrm{C}$，峰值：$1550.9220\text{ nm},-20.61\text{ dBm}$；3dB衰减宽度为$0.180\text{ nm}$]{\includegraphics[width=\linewidth]{images/DFB-spectrum-20.jpg}}
    \caption{DFB 半导体激光器的光谱}
    \label{fig:DFB-spectrum}
\end{figure}
\begin{figure}[H]
    \centering
    \includegraphics[width=\linewidth]{images/FP-spectrum-14.jpg}
    \vspace{-2em}
    \caption{FP 半导体激光器的光谱，中心波长：1310.9360 nm，峰值：-7.75 dBm，平均纵模间隔：1.4270 nm}
    \label{fig:FP-spectrum}
\end{figure}

\onecolumn
\section{预习报告}
\includepdf[pages=-]{images/预习报告-半导体激光器特性测量.pdf}

\end{document}
